\documentclass[11pt]{article}
\usepackage{ideastyle}

\begin{document}

\ideatitle{HillSmart}{Routes that understand hills, stairs, and effort, with hands-free voice guidance.}

\section{The Idea}
Ithaca is famous for its hills. Normal maps pick the shortest route, even if it
climbs Libe Slope or uses stairs. HillSmart lets you choose what matters:
\emph{least climbing}, \emph{no stairs} (wheelchair, stroller, scooter, bike), or
\emph{fastest}. It shows the elevation profile of each option and guides you
hands-free by voice, which is useful on a bike or with a wheelchair.

\section{Track Fit}
\begin{tabular}{@{}P{0.28\textwidth}P{0.66\textwidth}@{}}
\toprule
\req{Built with Cursor}{Routing, elevation analysis, and UI built in Cursor.}
\req{Grok Voice / Imagine}{\textbf{Voice} turn-by-turn and questions (\emph{``Is there a flatter way?''}).}
\req{Space data}{Not needed. (Optional twist: elevation data comes from satellite measurements.)}
\req{Navigation theme}{Direct fit; accessibility angle.}
\req{Also eligible}{Big Red, Software, Design.}
\bottomrule
\end{tabular}

\section{Features}
\subsection{MVP}
\begin{itemize}
  \item Enter start and destination (by voice or text).
  \item Get 2--3 candidate routes; compute total climb and steepest section for each.
  \item Elevation profile chart and a map with the route color-coded by steepness.
  \item Voice summary: \emph{``Route B adds 3 minutes but saves 40 meters of climbing.''}
\end{itemize}
\subsection{Stretch}
\begin{itemize}
  \item Custom routing that truly avoids steep segments (not just picking between options).
  \item Crowd-reported obstacles (broken elevator, icy path).
  \item Live voice turn-by-turn while moving.
\end{itemize}

\section{Tech Stack and Data}
\begin{itemize}
  \item Routes: Google Maps Directions API (alternatives) or OpenStreetMap with OSRM/GraphHopper.
  \item Elevation: Google Elevation API or open elevation data.
  \item Frontend: React with Leaflet/Google Maps; a chart library for profiles.
  \item AI: Grok Voice API.
\end{itemize}

\section{Limitations and Risks}
\begin{itemize}
\limitation{Picking vs.\ planning}{Choosing among Google's alternative routes is easy but limited; real slope-aware routing needs your own routing engine.}{Ship ``compare alternatives'' first; custom routing is a stretch goal.}
\limitation{Stairs data is incomplete}{OpenStreetMap marks many but not all stairs and ramps on campus.}{Hand-check the demo routes and add crowd reports.}
\limitation{Elevation resolution}{Public elevation data can miss short steep sections and stairs.}{Treat numbers as estimates; combine with stairs tags.}
\limitation{Existing features}{Google Maps shows wheelchair-accessible transit and bike elevation in some places.}{Focus on walking/wheelchair hill avoidance on campus, which Google doesn't do well.}
\limitation{API cost and keys}{Google Maps APIs need billing set up.}{Use free credits or OpenStreetMap; never commit the key.}
\end{itemize}

\section{Scorecard}
\begin{tabular}{@{}P{0.25\textwidth}P{0.08\textwidth}P{0.6\textwidth}@{}}
\toprule
\rating{Demo-able}{4}{Clear map demo; anyone at Cornell relates.}
\rating{Buildable in 24h}{4}{MVP is mostly APIs.}
\rating{Navigation theme}{5}{Direct fit.}
\rating{Wow factor}{3}{Useful, not flashy.}
\rating{Technical risk (5 = riskiest)}{2}{Low (MVP), moderate (custom routing).}
\bottomrule
\end{tabular}

\section{Verdict}
\textbf{Good practical pick.} Very relatable to Cornell judges. The voice part
must feel essential, not added on.

\end{document}
